\begin{solapartat}
El flux creat per un camp magnètic en una espira ve determinat per:
\[ \Phi = \vec{B}\cdot\vec{S} = B\,S\,\cos(\alpha) \quad \text{\punts{0,3}} \]
on $\alpha$ és l'angle que forma la direcció del camp magnètic amb la perpendicular a l'espira, per tant $\alpha = 60°$
\[ \Phi = 0,5\,\pi\,0,04^2\,\cos(60°) = 1,3\cdot 10^{-3}\ \mathrm{Wb} \quad \text{\punts{0,4}} \]
Donat que el flux que travessa l'espira es constant en el temps, no s'induirà cap \textit{fem}. \punts{0,3}
\end{solapartat}

\begin{solapartat}
Per la gràfica que ens mostren en el enunciat el camp magnètic varia linealment segons l'expressió:
\[ B(t) = 0,5 - \frac{0,5}{100\cdot 10^{-3}}\,t \quad \text{\punts{0,3}} \]
Per tant el flux que genera el camp serà:
\[ \Phi(t) = \pi\,0,04^2\,\left(0,5 - \frac{0,5}{100\cdot 10^{-3}}\,t\right)\,\cos(60°) \quad \text{\punts{0,3}} \]
i la \textit{fem} generada serà:
\[ \varepsilon = -\frac{\mathrm{d}\Phi}{\mathrm{d}t} = \pi\,0,04^2\,\cos(60°)\,\frac{0,5}{100\cdot 10^{-3}} = 1,3\cdot 10^{-2}\ \mathrm{V} \quad \text{\punts{0,4}} \]
\end{solapartat}
